\documentclass[11pt,a4paper]{article}

\usepackage[T1]{fontenc}
\usepackage[utf8]{inputenc}
\usepackage[polish]{babel}
\usepackage{amsmath,amssymb,amsthm,mathtools}
\usepackage{geometry}
\usepackage{hyperref}
\usepackage{xcolor}
\usepackage{listings}
\usepackage{booktabs}
\usepackage{array}
\usepackage{enumitem}
\usepackage{microtype}

\geometry{margin=2.4cm}

% ---------- skróty matematyczne ----------
\newcommand{\C}{\mathbb{C}}
\newcommand{\R}{\mathbb{R}}
\newcommand{\Z}{\mathbb{Z}}
\newcommand{\F}{\mathbb{F}}
\newcommand{\id}{\mathbb{I}}
\newcommand{\1}{\mathbf{1}}
\newcommand{\bigO}{\mathcal{O}}
\newcommand{\diag}{\operatorname{diag}}
\newcommand{\popcount}{\operatorname{popcount}}
\newcommand{\supp}{\operatorname{supp}}
\newcommand{\sgn}{\operatorname{sgn}}
\newcommand{\ket}[1]{\left|#1\right\rangle}
\newcommand{\bra}[1]{\left\langle#1\right|}
\newcommand{\braket}[2]{\left\langle#1\,\middle|\,#2\right\rangle}
\newcommand{\w}{\omega}
\newcommand{\Alphabet}{\mathcal{A}}
\newcommand{\Cspace}{\mathcal{C}}
\newcommand{\code}[1]{\texttt{#1}}

% ---------- styl listingów ----------
\definecolor{codebg}{rgb}{0.97,0.97,0.97}
\definecolor{kw}{rgb}{0.10,0.10,0.55}
\definecolor{cmt}{rgb}{0.30,0.55,0.30}
\definecolor{str}{rgb}{0.55,0.10,0.10}

\lstdefinestyle{py}{%
  language=Python,
  basicstyle=\ttfamily\small,
  keywordstyle=\color{kw}\bfseries,
  commentstyle=\color{cmt}\itshape,
  stringstyle=\color{str},
  backgroundcolor=\color{codebg},
  showstringspaces=false,
  breaklines=true,
  numbers=left,
  numberstyle=\tiny\color{gray},
  numbersep=6pt,
  frame=single,
  framesep=3pt,
  framerule=0.2pt,
  rulecolor=\color{gray!40},
  xleftmargin=12pt,
  xrightmargin=2pt,
  literate=%
    {ą}{{\k a}}1 {ć}{{\'c}}1 {ę}{{\k e}}1 {ł}{{\l{}}}1
    {ń}{{\'n}}1 {ó}{{\'o}}1 {ś}{{\'s}}1 {ź}{{\'z}}1 {ż}{{\.z}}1
    {Ą}{{\k A}}1 {Ć}{{\'C}}1 {Ę}{{\k E}}1 {Ł}{{\L{}}}1
    {Ń}{{\'N}}1 {Ó}{{\'O}}1 {Ś}{{\'S}}1 {Ź}{{\'Z}}1 {Ż}{{\.Z}}1,
}
\lstset{style=py}

\hypersetup{%
  colorlinks=true,
  linkcolor=blue!50!black,
  urlcolor=blue!60!black,
  citecolor=blue!50!black,
  pdftitle={d6\_diagonal\_optimizer.py --- dokumentacja po polsku},
}

\title{\texttt{d6\_diagonal\_optimizer.py}\\[2pt]
       \large Optymalizacja diagonalnych bramek na dwóch ququdytach $d=6$\\
       w kodowaniu binarnym na 6 kubitach}
\author{Dokumentacja techniczna (po polsku, ze szczegółowymi wyjaśnieniami)}
\date{}

\begin{document}
\maketitle

\begin{abstract}
\noindent
Niniejszy dokument jest pełnym, polskojęzycznym opisem modułu
\code{d6\_diagonal\_optimizer.py}. Moduł syntetyzuje 6\nobreakdash-kubitową
realizację diagonalnej bramki $D[\Lambda]$ działającej na parze
ququdytów wymiaru $d=6$. Wykorzystuje fakt, że spośród $2^{6}=64$ stanów
bazowych przestrzeni Hilberta tylko $36$ leży w kodowej podprzestrzeni
qudytowej, a pozostałe $28$ to stany \emph{nielegalne} (\emph{don't-care}),
których fazy można wybrać dowolnie. Te 28 swobodnych faz jest
optymalizowanych po to, by powstała sieć bramek \code{CNOT}+\code{RZ}
była możliwie tania. Dokument tłumaczy każdą część modułu od podstaw:
matematykę transformaty Walsha, konwencje kodowania i indeksowania,
strukturę pipeline'u optymalizacyjnego (Random Search, Simulated
Annealing, relaksacja LP, Greedy, 2\nobreakdash-opt, polerowanie
\code{transpile}), pełne API wraz z opisem wszystkich argumentów oraz
wyniki empiryczne na dwóch przykładowych wektorach.
\end{abstract}

\tableofcontents
\bigskip\hrule\bigskip

% =====================================================================
\section{Sformułowanie problemu}
\label{sec:problem}
% =====================================================================

\subsection{Bramka diagonalna na dwóch ququdytach $d=6$}

Chcemy zaimplementować unitarną bramkę diagonalną
\begin{equation}
\label{eq:Dlambda}
D[\Lambda] \;=\; \diag\!\bigl(e^{i\theta_{0}},\,e^{i\theta_{1}},\,\dots,\,
                              e^{i\theta_{35}}\bigr)
\end{equation}
działającą na dwóch qudytach wymiaru $d=6$. Operator $D[\Lambda]$ jest
$36\times 36$, wstępnie określony przez 36 rzeczywistych faz
$\theta_{0},\dots,\theta_{35}$. Każdy poziom qudytowy
$a\in\{0,1,2,3,4,5\}$ jest kodowany w trzech kubitach fizycznych za pomocą
zwykłego kodu binarnego:
\begin{equation}
\label{eq:encoding}
\begin{aligned}
0 &\;\mapsto\; \ket{000}, &
1 &\;\mapsto\; \ket{001}, &
2 &\;\mapsto\; \ket{010}, \\[1pt]
3 &\;\mapsto\; \ket{011}, &
4 &\;\mapsto\; \ket{100}, &
5 &\;\mapsto\; \ket{101}.
\end{aligned}
\end{equation}
Dwa stany $\ket{110}$ oraz $\ket{111}$ każdego trójkubitowego bloku
\emph{nie} należą do przestrzeni kodowej. Łącznie spośród $2^{6}=64$
stanów bazowych przestrzeni 6\nobreakdash-kubitowej $36$ jest
\emph{legalnych} (te, które kodują parę $(a,b)\in\{0,\dots,5\}^{2}$),
a $28$ jest \emph{nielegalnych} (\emph{don't-care}).

\subsection{Stopnie swobody i cel optymalizacji}

Bramka $D[\Lambda]$ musi działać poprawnie tylko na 36 stanach
legalnych. Na 28 stanach nielegalnych może działać dowolnie -- w
szczególności wolno jej tam mieć \emph{dowolne} fazy. Dobrze dobrane
fazy nielegalne potrafią drastycznie uprościć obwód (zob.\
sekcja~\ref{sec:why-helps}).

Zadanie modułu rozkłada się na cztery kroki:
\begin{enumerate}[leftmargin=*]
  \item dobrać 28 faz \emph{nielegalnych}, traktując je jako parametry
        rzeczywiste,
  \item zsyntetyzować unitarną diagonalną dla powstałego wektora długości
        64 jako sieć \emph{phase gadgets} (\code{CNOT}+\code{RZ}),
  \item porównać uzyskany obwód z baseline'em z biblioteki Qiskit
        (\code{DiagonalGate}),
  \item zwrócić oba obwody wraz z ich kosztami zasobowymi (głębokość,
        liczba \code{CX}, liczba \code{RZ}).
\end{enumerate}

% =====================================================================
\section{Konwencje}
\label{sec:conventions}
% =====================================================================

Wszystkie konwencje poniżej są ustalone i konsekwentnie stosowane w
całym module.

\subsection{Endianness i indeks bazy}
\label{sec:endianness}

Indeks stanu bazowego pary $(a,b)$ definiujemy jako
\begin{equation}
\label{eq:index}
\mathrm{index}(a,b) \;=\; (a \ll 3) \,\lor\, b \;=\; 8a+b,
\qquad a,b\in\{0,\dots,5\}.
\end{equation}
Tutaj $\ll$ to przesunięcie bitowe w lewo, a $\lor$ to bitowe \emph{or}
(równoważnie dodawanie, bo bity nie kolidują). Oznacza to, że:
\begin{itemize}[leftmargin=*]
  \item kubity numerowane jako \textbf{5, 4, 3} przenoszą trzybitowy
        kod binarny qudytu $a$ (kubit~5 jest najbardziej znaczący),
  \item kubity numerowane jako \textbf{2, 1, 0} przenoszą trzybitowy
        kod binarny qudytu $b$.
\end{itemize}
W konwencji Qiskit-a indeks $k\in\{0,\dots,63\}$ stanu bazowego
odpowiada konfiguracji, w której fizyczny kubit $i$ przyjmuje wartość
$(k \gg i)\,\&\,1$ (gdzie $\gg$ to przesunięcie w prawo, a~$\&$ to bitowe
\emph{and}). Innymi słowy, najmłodszy bit $k$ to wartość kubitu~0.

\paragraph{Indeksy legalne i nielegalne.} Zbiór indeksów legalnych to
dokładnie obrazy~\eqref{eq:index} dla $a,b\in\{0,\dots,5\}$:
\begin{equation}
\label{eq:legal}
L \;=\; \bigl\{\,8a+b\;:\;0\le a,b\le 5\,\bigr\},\qquad |L|=36.
\end{equation}
Pozostałe $28$ indeksów stanowi $L^{c}\subset\{0,\dots,63\}$:
\begin{lstlisting}
from d6_diagonal_optimizer import legal_indices, illegal_indices
assert len(legal_indices()) == 36
assert len(illegal_indices()) == 28
assert set(legal_indices()).isdisjoint(illegal_indices())
\end{lstlisting}

\subsection{Kolejność wektora \texttt{diag36}}
\label{sec:order}

\code{diag36} jest płaskim wektorem długości 36 zawierającym fazy faz
legalnych. Moduł obsługuje dwa porządki:
\begin{itemize}[leftmargin=*]
  \item \textbf{\code{order="row-major"}} (domyślnie): \\
        \code{diag36[6*a + b]} jest fazą stanu $\ket{a,b}$. Odpowiada
        to czytaniu tabeli $6\times6$ wierszami (indeksowanej przez
        $a$ w wierszach, $b$ w kolumnach), zgodnie z notacją artykułu
        AME(4,6).
  \item \textbf{\code{order="col-major"}}: \\
        \code{diag36[a + 6*b]} jest fazą stanu $\ket{a,b}$.
\end{itemize}
Wybór tej opcji wpływa wyłącznie na to, w jaki sposób użytkownik
indeksuje wejście; wszystkie obliczenia wewnętrzne korzystają z
indeksu~\eqref{eq:index}.

\subsection{Format wejściowy}

\code{diag36} może być podany jako:
\begin{itemize}[leftmargin=*]
  \item tablica długości 36 \textbf{liczb zespolonych} o module 1
        (np.\ \code{[1, omega, omega.conjugate(), \dots]}); moduł
        wyciąga z nich kąt fazowy przez \code{np.angle}, \emph{albo}
  \item tablica długości 36 \textbf{rzeczywistych faz w radianach}.
\end{itemize}

\paragraph{Auto-detekcja typu (funkcja \code{\_phases\_from\_diag36}).}
Algorytm rozpoznawania typu wejścia działa następująco:
\begin{enumerate}[leftmargin=*]
  \item jeżeli \code{dtype} jest zespolony lub którakolwiek z części
        urojonych jest niezerowa $\Rightarrow$ traktujemy wejście jako
        zespolone i bierzemy $\theta_{j}=\arg(\code{diag36}[j])$;
  \item w przeciwnym razie, jeśli wszystkie wpisy mają wartość
        bezwzględną równą $1$ z tolerancją numeryczną
        $\Rightarrow$ traktujemy wejście jako rzeczywiste $\pm 1$
        i ponownie bierzemy fazę;
  \item w pozostałych przypadkach $\Rightarrow$ traktujemy wpisy jako
        \emph{surowe kąty} fazowe w radianach.
\end{enumerate}

\subsection{Normalizacja faz}
\label{sec:normalize}

Przed transformatą Walsha wszystkie fazy są \emph{zawijane} do
półotwartego przedziału $(-\pi,\pi]$ funkcją \code{normalize\_phase}.
Każde \emph{podniesienie} $\theta+2\pi k(x)$ dla całkowitoliczbowego
$k(x)$ nie zmienia $e^{i\theta(x)}$ -- jest to swobodny wybór, ale
różny wybór $k(\cdot)$ daje różne współczynniki Walsha (a więc różne
syntezy). Moduł zawsze stosuje \emph{kanoniczną} reprezentację z
wartością główną.

\subsection{Konwencja \texttt{RZ} w Qiskicie}
\label{sec:rz-convention}

W Qiskicie obowiązuje
\begin{equation}
\label{eq:rz}
\mathrm{RZ}(\varphi) \;=\; e^{-\,i\,\varphi\,Z/2}
\;=\; \diag\!\bigl(e^{-i\varphi/2},\,e^{i\varphi/2}\bigr).
\end{equation}
Aby $e^{i\,c\,Z}$ pokrywało się z~\eqref{eq:rz}, trzeba mieć
$\varphi=-2c$. Wielokubitowy operator
$\exp(i\,c\,Z_{S})$, gdzie $Z_{S}=\bigotimes_{i\in S}Z_{i}$, jest
realizowany jako:
\begin{enumerate}[leftmargin=*]
  \item drabinka \code{CNOT} przenosząca parzystość bitów na ostatni
        kubit nośnika $S$ (\emph{target}),
  \item pojedyncza bramka $\mathrm{RZ}(-2c)$ na tym kubicie,
  \item odwrócona drabinka \code{CNOT} odkomputowująca parzystość.
\end{enumerate}
Dla $|S|=1$ drabinka jest pusta i emitujemy tylko \code{RZ}. Zgodność
ze \code{qiskit.circuit.library.DiagonalGate} jest weryfikowana
numerycznie w \code{\_self\_tests}.

% =====================================================================
\section{Rozwinięcie Walsha (wielomian fazowy)}
\label{sec:walsh}
% =====================================================================

\subsection{Definicja i interpretacja}

Dla unitarnej diagonalnej $D=\diag\bigl(e^{i\theta_{0}},\dots,
e^{i\theta_{N-1}}\bigr)$ na $n$ kubitach, $N=2^{n}$, każdą fazę
zapisujemy w bazie znaków grupy $\F_{2}^{\,n}$:
\begin{equation}
\label{eq:walsh-expansion}
\theta(x) \;=\; \sum_{S\in\{0,1\}^{n}} c_{S}\,(-1)^{\popcount(S \,\&\, x)},
\qquad x\in\{0,\dots,N-1\}.
\end{equation}
Tutaj $S\,\&\, x$ to bitowe \emph{and}, a $\popcount(\cdot)$ to liczba
jedynek w zapisie binarnym (\emph{Hamming weight}). Współczynniki
$c_{S}\in\R$ to transformata Walsha--Hadamarda funkcji $\theta(\cdot)$:
\begin{equation}
\label{eq:walsh-coeff}
c_{S} \;=\; \frac{1}{N}\sum_{x=0}^{N-1}\theta(x)\,(-1)^{\popcount(S\,\&\, x)}.
\end{equation}
Równoważnie:
\begin{equation}
\label{eq:walsh-D}
D \;=\; e^{i c_{0}}\cdot\id\,\cdot\;\prod_{S\neq 0}\,e^{\,i\,c_{S}\,Z^{S}},
\qquad
Z^{S}\;=\;\bigotimes_{i\in S} Z_{i}.
\end{equation}
Każdy niezerowy współczynnik $c_{S}$ odpowiada jednemu \emph{phase
gadgetowi} (sekcja~\ref{sec:rz-convention}): drabince \code{CNOT}
liczącej parzystość $\bigoplus_{i\in S} x_{i}$, jednej rotacji
\code{RZ($-2 c_{S}$)} na targecie i odwróconej drabince. Wyraz $c_{0}$
to globalna faza -- w obwodzie zwijamy ją do \code{qc.global\_phase}
i nigdy nie syntetyzujemy.

\subsection{Implementacja transformaty}

Transformata Walsha jest liczona w miejscu \emph{butterfly}-em w
$\bigO(N\log N)$ ze znakową konwencją
$h[m]=\sum_{k}(-1)^{\popcount(m\&k)}\theta_{k}$, a następnie skalowana
przez $1/N$. Tolerancja numeryczna na obcięcie małych współczynników
(parametr \code{tol}) wynosi domyślnie $10^{-10}$. Test \emph{round-trip}
\code{theta -> walsh -> theta} jest wykonywany w
\code{\_self\_tests}.

% =====================================================================
\section{Dlaczego optymalizacja faz \emph{don't-care} pomaga?}
\label{sec:why-helps}
% =====================================================================

\subsection{Koszt syntezy zależy od liczby niezerowych $c_{S}$}

Zarówno synteza Qiskit-owa (\code{DiagonalGate}) jak i synteza poprzez
\emph{phase gadgety} płacą jedną bramkę \code{RZ} (oraz pewną liczbę
\code{CX}) za \emph{każdy} niezerowy współczynnik Walsha $c_{S}$.
Konkretnie:
\begin{itemize}[leftmargin=*]
  \item \textbf{Phase gadgets}: drabinka \code{CNOT} dla nośnika
        $S\subseteq\{0,\dots,n-1\}$ kosztuje $2(|S|-1)$ bramek
        \code{CX} (drabinka i jej odwrócenie) oraz jedną \code{RZ}.
  \item \textbf{\code{DiagonalGate}}: Qiskit używa konstrukcji typu
        Shende (rząd $2^{n}-2$ bramek \code{CX} w najgorszym razie),
        natomiast późniejszy \code{transpile} eliminuje wszystkie
        rotacje o kącie dokładnie zerowym, co często oszczędza
        bardzo wiele \code{CX}.
\end{itemize}
W obu przypadkach: \emph{im mniej niezerowych $c_{S}$, tym tańszy
obwód.}

\subsection{Afiniczna zależność od faz nielegalnych}

Fazy 36 indeksów \emph{legalnych} są ustalone przez $D[\Lambda]$.
Pozostałych 28 faz \emph{nielegalnych}
$\theta^{\mathrm{ill}}\in\R^{28}$
to swobodne parametry. Z definicji
transformaty~\eqref{eq:walsh-coeff} każdy współczynnik $c_{S}$ jest
\emph{afiniczną} funkcją $\theta^{\mathrm{ill}}$:
\begin{equation}
\label{eq:cS-affine}
c_{S} \;=\; \mathrm{const}_{S} \;+\; \sum_{i=0}^{27}
B_{S,i}\,\theta^{\mathrm{ill}}_{i},
\qquad
B_{S,i} \;=\; \frac{(-1)^{\popcount(S\,\&\, k^{\mathrm{ill}}_{i})}}{N},
\end{equation}
gdzie $k^{\mathrm{ill}}_{i}$ jest $i$-tym indeksem nielegalnym, a
$\mathrm{const}_{S}$ to wkład 36 indeksów legalnych. Macierz $B$ ma
rozmiar $63\times 28$ (bez wiersza $S=0$, bo ten jest globalną fazą),
a jej wpisy są równe $\pm 1/N$.

\paragraph{Przestrzeń poszukiwań.} Wektory $(c_{S})_{S\neq 0}$ leżą
w 28-wymiarowej afinicznej podprzestrzeni $\R^{63}$. Po ograniczeniu do
dyskretnego alfabetu (np.\ $\Alphabet=\{0,\,2\pi/3,\,-2\pi/3\}$) zostaje
\emph{krata} $|\Alphabet|^{28}=3^{28}\approx 2{,}3\times 10^{13}$
konfiguracji. Cel: znaleźć punkt o jak najmniejszej liczbie niezerowych
$c_{S}$ (równoważnie, najrzadszy punkt w afinicznej podprzestrzeni).

\subsection{Trudność problemu}

To w ogólności trudny problem kombinatoryczny / odzyskiwania rzadkiego
sygnału ($\ell_{0}$-minimalizacji). Moduł atakuje go \emph{stosem
heurystyk}: random search $\to$ multi-restart simulated annealing $\to$
relaksacja LP $\to$ greedy coordinate descent $\to$ pairwise 2-opt,
i opcjonalnie poleruje wynik metryką, którą obliczamy faktycznym
\code{transpile}'em (sekcja~\ref{sec:pipeline}).

% =====================================================================
\section{Pipeline optymalizatora}
\label{sec:pipeline}
% =====================================================================

W trybie domyślnym \code{method="best\_of"} etapy są wykonywane w
następującej kolejności i każdy aktualizuje wspólny rekord
$(\theta^{\mathrm{ill}}_{\mathrm{best}}, \mathrm{cost}_{\mathrm{best}})$:

\begin{equation*}
\boxed{\;
\begin{array}{l}
\text{(1) random\_search}\\
\;\;\downarrow\\
\text{(2) multi\_restart\_simulated\_annealing}\;\;(n_\mathrm{restarts}=6)\\
\;\;\downarrow\\
\text{(3) lp\_relax\_walsh\_l1}\;+\;\text{projekcja na alfabet}\\
\;\;\downarrow\\
\text{(4) lp(rounded)}\;\to\;\text{greedy\_coordinate\_descent}\\
\;\;\downarrow\\
\text{(5) multi\_restart\_greedy}\;\;(n_\mathrm{restarts}=6)\\
\;\;\downarrow\\
\text{(6) two\_opt\_descent}\;\;(\text{finalny pairwise local search})\\
\;\;\downarrow\\
\text{(7) opcjonalny polish\_target z transpile()}
\end{array}\;}
\end{equation*}

Każdy etap dostaje \emph{cały} budżet \code{max\_iters}; dodawanie
nowych etapów może co najwyżej wzmocnić wyszukiwanie (kosztem czasu) --
nigdy go nie osłabi. Poniżej szczegółowy opis każdego etapu.

% ---------------------------------------------------------------------
\subsection{Random search (RS)}
\label{sec:rs}

Funkcja \code{random\_search(legal\_phases, ..., alphabet, max\_iters,
seed)}. Algorytm:
\begin{enumerate}[leftmargin=*]
  \item w każdej iteracji losujemy 28 faz nielegalnych \emph{i.i.d.}
        z dyskretnego alfabetu $\Alphabet$;
  \item budujemy pełen wektor $\theta\in\R^{64}$ przez wstawienie 36
        legalnych faz na pozycje legalne i 28 wylosowanych na pozycje
        nielegalne (\code{\_build\_theta64});
  \item liczymy \code{phase\_polynomial\_cost} (zob.\
        sekcja~\ref{sec:cost});
  \item zapamiętujemy konfigurację o najmniejszym koszcie.
\end{enumerate}
Tania, deterministyczna (przy ustalonym \code{seed}) baseline'owa
strategia. Dla naturalnego alfabetu
$\Alphabet=\{0,\,2\pi/3,\,-2\pi/3\}$ przestrzeń ma $3^{28}\approx
2{,}3\times 10^{13}$ punktów -- nawet 4000 losowań wyłapuje rozsądne
\emph{baseny przyciągania} dla późniejszych etapów.

% ---------------------------------------------------------------------
\subsection{Multi-restart simulated annealing}
\label{sec:sa}

Funkcja \code{multi\_restart\_simulated\_annealing(...)} uruchamia
\code{simulated\_annealing} $n_{\mathrm{restarts}}$ razy:
\begin{itemize}[leftmargin=*]
  \item \emph{pierwszy} restart startuje \emph{na ciepło} z dotychczas
        najlepszej konfiguracji (typowo z RS),
  \item kolejne restarty startują z nowych losowych konfiguracji.
\end{itemize}
Każdy pojedynczy SA wykonuje ruchy \emph{pojedynczych zmian} (single-flip
Metropolis): w każdej iteracji wybieramy losowy indeks $i\in\{0,\dots,27\}$
i losową wartość $a'\in\Alphabet$; akceptujemy zmianę
$\theta^{\mathrm{ill}}_{i}\leftarrow a'$ z prawdopodobieństwem
$\min(1, e^{-\Delta/T})$, gdzie $\Delta$ to przyrost kosztu, a $T$ to
aktualna temperatura. Schemat chłodzenia jest geometryczny od
$T_{\mathrm{start}}=\max(2,\,0{,}5\cdot\mathrm{cost})$ do
$T_{\mathrm{end}}=0{,}05$.

\paragraph{Po co restarty.} Restarty są głównym lekarstwem na
tendencję pojedynczego SA do utykania w tym samym lokalnym minimum,
w którym leży RS-warmstart. Sześć restartów jest tanim, dobrze
zbadanym kompromisem.

% ---------------------------------------------------------------------
\subsection{Relaksacja LP $\min \|c\|_{1}$}
\label{sec:lp}

Funkcja \code{lp\_relax\_walsh\_l1(legal\_phases, legal\_idx,
illegal\_idx)} rozwiązuje \emph{ciągłą} wypukłą relaksację
\begin{equation}
\label{eq:lp}
\min_{\theta^{\mathrm{ill}}\in[-\pi,\pi]^{28}}\;
\sum_{S\neq 0}\;\bigl|\,\mathrm{const}_{S}+B_{S,:}\,\theta^{\mathrm{ill}}\,\bigr|.
\end{equation}
Po standardowej linearyzacji modułu wprowadzamy zmienne pomocnicze
$u_{S}\ge 0$ takie, że
$-u_{S}\le\mathrm{const}_{S}+B_{S,:}\theta^{\mathrm{ill}}\le u_{S}$.
W rezultacie mamy LP w
\begin{itemize}[leftmargin=*]
  \item $28+63=91$ zmiennych ($\theta^{\mathrm{ill}}$ + zmienne $u$),
  \item $2\cdot 63=126$ nierównościach ograniczających moduł,
  \item plus pudełkowe ograniczenia $[-\pi,\pi]$ na $\theta^{\mathrm{ill}}$.
\end{itemize}
Rozwiązanie: solver HiGHS przez \code{scipy.optimize.linprog}.

\paragraph{Dlaczego $\ell_{1}$.} Minimalizacja $\ell_{1}$ jest
standardową \emph{wypukłą} relaksacją problemu $\ell_{0}$ (liczby
niezerowych) i sprzyja rzadkim rozwiązaniom (zob.\ teorię
\emph{compressed sensing}). Po znalezieniu rozwiązania ciągłego
projektujemy je per-koordynata na najbliższy punkt alfabetu (funkcja
\code{\_project\_to\_alphabet}, dystans cykliczny
$\bmod 2\pi$), a następnie polerujemy \code{greedy\_coordinate\_descent}.

\paragraph{Praktyka.} Dla wektorów z sekcji~\ref{sec:results}
\emph{zaokrąglony} punkt LP nie jest dyskretnym optimum -- legalne fazy
żyją w podkratownicy \emph{trzecich pierwiastków z 1}, a LP nie widzi
struktury dyskretnej. Niemniej dostarcza \emph{innego} startu, co bywa
istotne, gdy etapy losowe i SA skleiły się w jednym basenie.

% ---------------------------------------------------------------------
\subsection{Greedy coordinate descent (1-opt)}
\label{sec:greedy}

Funkcja \code{greedy\_coordinate\_descent(legal\_phases, ...,
init\_phases, eval\_cost, max\_passes=12)}. Algorytm:
\begin{enumerate}[leftmargin=*]
  \item rozpoczynamy z \code{init\_phases};
  \item w każdym \emph{przebiegu} chodzimy po 28 indeksach nielegalnych
        kolejno; dla każdego sprawdzamy \emph{wszystkie} wartości
        alfabetu i przyjmujemy tę o najmniejszym koszcie;
  \item kończymy, gdy w pełnym przebiegu żaden indeks nie poprawił
        kosztu (lub po \code{max\_passes} przebiegach).
\end{enumerate}
Przy $|\Alphabet|=3$ i 28 indeksach pojedynczy przebieg wymaga
$28\cdot 3=84$ ewaluacji kosztu; konwergencja wymaga zazwyczaj 3--5
przebiegów. Tani, w pełni deterministyczny, niezawodnie schodzi do
lokalnego minimum (1-opt).

Funkcja \code{multi\_restart\_greedy(...)} uruchamia 6 \emph{greedy}-ów
z losowych startów (i opcjonalnie jeden \emph{warm-start} z bieżącego
najlepszego), zachowując konfigurację o najmniejszym koszcie.

% ---------------------------------------------------------------------
\subsection{Pairwise 2-opt local search}
\label{sec:2opt}

Funkcja \code{two\_opt\_descent(legal\_phases, ..., init\_phases,
eval\_cost, max\_passes=3)}. Algorytm:
\begin{enumerate}[leftmargin=*]
  \item dla każdej \emph{nieuporządkowanej} pary $(i,j)$ indeksów
        nielegalnych, $i<j$;
  \item dla każdego \emph{wspólnego} przypisania
        $(a_{i}, a_{j})\in\Alphabet\times\Alphabet$;
  \item przyjmujemy parę o najmniejszym koszcie (jeśli poprawia
        bieżący stan).
\end{enumerate}
Koszt pojedynczego przebiegu przy $|\Alphabet|=3$:
$\binom{28}{2}\cdot 9 = 378\cdot 9 = 3402$ ewaluacji.

\paragraph{Dlaczego 2-opt jest ważny.} W eksperymentach to ten etap
\emph{przesuwa najwięcej} (na wektorze \code{D\_Lambda\_23}: liczba
wyrazów Walsha $45\to 33$, głębokość baseline'u $94\to 80$). Greedy
1-opt schodzi do takiego minimum, w którym żadna \emph{pojedyncza}
zmiana nie poprawi kosztu, ale wiele takich minimów daje się jeszcze
opuścić zmieniając \emph{dwa} indeksy jednocześnie. 2-opt pokrywa
dokładnie te kierunki. 3-opt i wyższe nie są zaimplementowane, bo
sześcienny narzut rzadko się opłaca.

% ---------------------------------------------------------------------
\subsection{Opcjonalne polerowanie z \texttt{transpile}}
\label{sec:polish}

Parametr \code{polish\_target} (domyślnie \code{"cx\_then\_depth"})
włącza ostatnią rundę \code{greedy\_coordinate\_descent}, w której
funkcja kosztu jest wyliczana \emph{realnym} \code{transpile}'em:
\begin{lstlisting}
def polish_eval(theta64):
    return transpile_metric(
        theta64,
        target=polish_target,           # "depth" / "cx" / "weighted" / "cx_then_depth"
        via=polish_via,                 # "diagonal_gate" lub "phase_gadget"
        basis_gates=qiskit_basis,
        optimization_level=optimization_level,
        seed_transpiler=seed,
    )
\end{lstlisting}

\paragraph{Cele polerowania (\code{polish\_target}).}
\begin{itemize}[leftmargin=*]
  \item \code{"depth"} -- minimalizujemy głębokość obwodu po
        \code{transpile},
  \item \code{"cx"} -- minimalizujemy liczbę bramek \code{CX},
  \item \code{"weighted"} -- ważona kombinacja \code{cx} i głębokości,
  \item \code{"cx\_then\_depth"} -- leksykograficznie:
        najpierw \code{cx}, w razie remisu głębokość.
\end{itemize}
Polerowanie jest wolne (jedno \code{transpile} na ewaluację, $\sim$50
ms), ale optymalizuje \emph{dokładnie} tę metrykę, na której zależy
użytkownikowi. Na obu przykładowych wektorach z sekcji~\ref{sec:results}
polerowanie nie znajduje poprawy poza minimum kosztu Walshowego --
zastępczy koszt Walsha jest wystarczająco \emph{ciasny}. Polerowanie
domyślnie pozostaje włączone jako siatka bezpieczeństwa dla
przypadków, w których zastępczy koszt jest luźny. Można je wyłączyć
ustawiając \code{polish\_target=None}.

% ---------------------------------------------------------------------
\subsection{Funkcja kosztu}
\label{sec:cost}

\code{phase\_polynomial\_cost(theta64, ...)} zwraca
\code{CostBreakdown} o czterech polach skalarnych:

\begin{center}
\begin{tabular}{>{\ttfamily}l p{0.62\linewidth}}
\toprule
\rmfamily\textbf{Składnik} & \textbf{Znaczenie}\\
\midrule
num\_nonzero\_terms & liczba niezerowych współczynników Walsha (bez $S=0$)\\
cx\_estimate        & naiwny koszt \code{CX} dla phase gadgets:
                       $\displaystyle \sum_{S\,:\,c_{S}\neq 0} 2(\popcount(S)-1)$\\
weighted\_terms     & $\displaystyle \sum_{S\,:\,c_{S}\neq 0} \popcount(S)$\\
max\_support        & $\displaystyle \max_{S\,:\,c_{S}\neq 0} \popcount(S)$\\
\bottomrule
\end{tabular}
\end{center}

\noindent
Łączny skalar (domyślne wagi):
\begin{equation}
\label{eq:cost}
\mathrm{cost} \;=\; 1{,}0\cdot\code{cx\_estimate}
\;+\; 0{,}1\cdot\code{num\_nonzero\_terms}
\;+\; 0{,}05\cdot\code{max\_support}.
\end{equation}
Wagi można nadpisać argumentem
\code{cost\_weights=\{"cx":..., "terms":..., "support":..., "weighted":...\}}.
Optymalizator minimalizuje \eqref{eq:cost}; etap polerujący opcjonalnie
minimalizuje metrykę pochodzącą z \code{transpile}.

% =====================================================================
\section{Synteza phase gadgetów}
\label{sec:gadgets}
% =====================================================================

\code{synthesize\_phase\_gadgets(theta64, tol=1e-10, num\_qubits=6)}
buduje obiekt \code{QuantumCircuit(6)} zawierający wyłącznie bramki
\code{CX} oraz \code{RZ}, plus globalną fazę. Algorytm:
\begin{enumerate}[leftmargin=*]
  \item liczymy współczynniki Walsha $\{S\mapsto c_{S}\}$
        z tolerancją \code{tol};
  \item ustawiamy \code{qc.global\_phase = $c_{0}$} (tego wyrazu nie
        syntetyzujemy);
  \item sortujemy pozostałe pary $(S,c_{S})$ po
        $(\popcount(S),S)$ -- to czyni listę ładnie czytelną;
  \item dla każdego $(S,c_{S})$ liczymy nośnik $\supp(S)=
        [q_{0},\dots,q_{k-1}]$, ustawiamy target $t=q_{k-1}$ i
        emitujemy:
        \begin{itemize}
          \item drabinkę \code{CX}: $\code{cx}(q_{0},t),
                \code{cx}(q_{1},t),\dots,\code{cx}(q_{k-2},t)$,
          \item \code{rz($-2 c_{S}$, $t$)},
          \item odwróconą drabinkę \code{CX}.
        \end{itemize}
        Dla $k=1$ emitujemy tylko \code{rz}.
\end{enumerate}
Wynik jest \emph{ściśle} unitarną diagonalną; jest to weryfikowane w
\code{\_self\_tests} i w funkcji \code{extract\_diagonal\_from\_circuit},
która rzuca wyjątkiem, jeśli unitarne nie jest diagonalne w obrębie
tolerancji.

% =====================================================================
\section{Publiczne API}
\label{sec:api}
% =====================================================================

Po imporcie modułu dostępne są następujące nazwy:

\begin{lstlisting}
from d6_diagonal_optimizer import (
    # kodowanie / indeksowanie
    encode_qudit_level, pair_to_qubit_index, qubit_index_to_pair,
    legal_indices, illegal_indices,
    # obsluga wektora faz
    diag36_to_theta64, normalize_phase,
    # wielomian fazowy / koszt
    walsh_coefficients, phase_polynomial_cost,
    # synteza
    synthesize_phase_gadgets,
    # walidacja
    verify_on_code_space, extract_diagonal_from_circuit,
    # bloki budulcowe optymalizatora
    random_search, simulated_annealing,
    multi_restart_simulated_annealing,
    greedy_coordinate_descent, multi_restart_greedy,
    two_opt_descent, lp_relax_walsh_l1, transpile_metric,
    # glowny entry point
    optimize_d6_diagonal,
)
\end{lstlisting}

% ---------------------------------------------------------------------
\subsection{\texttt{optimize\_d6\_diagonal} -- główny punkt wejścia}
\label{sec:optimize-api}

\begin{lstlisting}
optimize_d6_diagonal(
    diag36,
    phase_alphabet="default",      # "default" | None | iterowalne float-ow
    max_iters=4000,
    seed=12345,
    qiskit_basis=None,             # domyslnie ["rz","sx","x","cx"]
    optimization_level=3,
    order="row-major",             # lub "col-major"
    method="best_of",              # "random"|"sa"|"greedy"|"lp"|"both"|"best_of"
    cost_weights=None,
    tol=1e-10,
    verbose=False,
    n_restarts=6,
    polish_target="cx_then_depth", # None | "depth" | "cx" | "weighted" | "cx_then_depth"
    polish_via="diagonal_gate",    # lub "phase_gadget"
    polish_max_passes=6,
) -> dict
\end{lstlisting}

\paragraph{\texttt{phase\_alphabet}.}
\begin{itemize}[leftmargin=*]
  \item \code{"default"} $\to$ alfabet $\{0,\,2\pi/3,\,-2\pi/3\}$.
        Naturalny dla wejść o wartościach w $\w_{3}$ (np.\ AME(4,6)).
  \item \code{None} $\to$ \emph{unikalne} fazy występujące we wpisach
        legalnych (zaokrąglone do 12 cyfr po przecinku). Trzymamy się
        wtedy alfabetu wejścia.
  \item dowolny iterowalny obiekt z liczbami zmiennoprzecinkowymi
        $\to$ używany dosłownie.
\end{itemize}

\paragraph{\texttt{method}.}
\begin{itemize}[leftmargin=*]
  \item \code{"random"} $\to$ tylko Random Search,
  \item \code{"sa"} $\to$ tylko (multi-restart) SA,
  \item \code{"greedy"} $\to$ tylko (multi-restart) Greedy,
  \item \code{"lp"} $\to$ LP $\to$ projekcja $\to$ greedy,
  \item \code{"both"} $\to$ RS $\to$ multi-restart SA (poprzedni
        domyślny),
  \item \code{"best\_of"} $\to$ pełen pipeline z sekcji~\ref{sec:pipeline}
        (obecny domyślny).
\end{itemize}

\paragraph{Wartość zwracana.}
Funkcja zwraca słownik z artefaktami:

\begin{center}
\begin{tabular}{>{\ttfamily}l p{0.62\linewidth}}
\toprule
\rmfamily\textbf{Klucz} & \textbf{Znaczenie}\\
\midrule
best\_theta64           & wektor faz długości 64 (legalne z \code{diag36},
                          nielegalne dobrane przez optymalizator)\\
best\_diag64            & \code{np.exp(1j * best\_theta64)}\\
best\_illegal\_phases   & lista długości 28 -- wybrane fazy \emph{don't-care}\\
legal\_indices,
illegal\_indices        & listy 36 i 28 indeksów\\
walsh\_coeffs           & słownik \code{\{mask: c\_S\}} po obcięciu tolerancją\\
num\_nonzero\_terms,
cx\_estimate,
max\_support            & podsumowujące liczby z \code{phase\_polynomial\_cost}\\
alphabet                & faktycznie użyty alfabet\\
method,
polish\_target,
polished                & co zostało uruchomione i czy polerowanie pomogło\\
phase\_gadget\_circuit  & naiwny obwód \code{CX}+\code{RZ} (przed transpile)\\
transpiled\_circuit     & ten sam po \code{transpile} do żądanej bazy\\
baseline\_circuit       & Qiskitowy \code{DiagonalGate(64)} -> \code{transpile},
                          z fazami \emph{don't-care} dobranymi przez optymalizator\\
metrics                 & głębokość/\code{cx}/\code{rz} oraz rozkład
                          per-bramka dla wszystkich trzech obwodów\\
\bottomrule
\end{tabular}
\end{center}

% ---------------------------------------------------------------------
\subsection{Funkcje walidujące}
\label{sec:validate}

\paragraph{\texttt{verify\_on\_code\_space(qc, diag36, order="row-major")}.}
Zwraca \emph{największy} z modułów odchyleń pomiędzy diagonalą obwodu na
36 indeksach legalnych a celem \code{diag36}, po wyrównaniu globalnej
fazy w punkcie referencyjnym o największej wartości bezwzględnej.
Indeksy \emph{don't-care} są ignorowane. Próg
$<10^{-7}$ jest stosowany jako rygorystyczna asercja w testach.

\paragraph{\texttt{extract\_diagonal\_from\_circuit(qc, atol=1e-9)}.}
Zwraca diagonalę macierzy \code{Operator(qc).data}. Rzuca
\code{ValueError}, jeśli unitarne nie jest diagonalne w obrębie
\code{atol}.

% ---------------------------------------------------------------------
\subsection{Self-testy}
\label{sec:tests}

Funkcja \code{\_self\_tests()} (uruchamiana z bloku \code{\_\_main\_\_})
sprawdza:
\begin{enumerate}[leftmargin=*]
  \item helpery kodowania/indeksów: zbiory legalnych i nielegalnych
        są rozłączne, ich suma to \code{range(64)}, każdy indeks
        legalny rozkłada się na $(a,b)$ z $a,b<6$;
  \item \emph{round-trip} transformaty Walsha--Hadamarda na losowym
        wektorze;
  \item zgodność \code{synthesize\_phase\_gadgets} z
        \code{qiskit.DiagonalGate} z dokładnością do globalnej fazy
        na losowej diagonali;
  \item \code{verify\_on\_code\_space} $<10^{-7}$ na losowym
        \code{diag36}.
\end{enumerate}

% =====================================================================
\section{Wyniki empiryczne}
\label{sec:results}
% =====================================================================

Wszystkie liczby poniżej są dla parametrów domyślnych:
\begin{itemize}[leftmargin=*]
  \item \code{phase\_alphabet="default"},
  \item \code{max\_iters=4000}, \code{seed=2024},
  \item \code{method="best\_of"}, \code{polish\_target="cx\_then\_depth"},
  \item \code{basis\_gates=["rz","sx","x","cx"]},
        \code{optimization\_level=3},
\end{itemize}
na Qiskit 2.3.1.

% ---------------------------------------------------------------------
\subsection{Wektor A -- oryginalne $D[\Lambda_{2,3}]$
            (\texttt{\_example\_d\_lambda\_2\_3})}
\label{sec:vec-A}

\begin{lstlisting}
omega3     = np.exp(2j * np.pi / 3)
omega3_bar = np.conj(omega3)
diag = np.array([
     1, omega3, omega3_bar, omega3, omega3_bar, 1,
     1,      1,      omega3,      1, omega3_bar, 1,
     1, omega3,           1, omega3_bar,         1, 1,
     1, omega3_bar,       1, omega3, omega3, omega3,
     omega3_bar, omega3_bar, omega3_bar, omega3_bar, 1, 1,
     omega3, omega3, omega3_bar, 1, 1, omega3_bar,
])
\end{lstlisting}

\begin{center}
\begin{tabular}{l r r r r}
\toprule
\textbf{obwód} & \textbf{depth} & \textbf{cx} & \textbf{rz} &
\textbf{niez.\ wyr.\ Walsha}\\
\midrule
naiwny baseline (\emph{don't-care} = 1) & --- & --- & --- & ---\\
phase-gadget (raw)                       & 173 & 150 & 41 & 41\\
phase-gadget po \code{transpile}         & 125 & 102 & 41 & 41\\
\code{DiagonalGate} baseline
   (z optymalizowanymi \emph{don't-care}) & \textbf{85} & \textbf{50} & 41 & 41\\
\bottomrule
\end{tabular}
\end{center}

\noindent
\code{verify\_on\_code\_space} $\approx 1{,}5\times 10^{-15}$ dla
wszystkich trzech wariantów. Wszystkie etapy optymalizatora skupiły się
na tym samym koszcie Walshowym
(\code{cost}$=154{,}350$, \code{cx\_estimate}$=150$,
\code{terms}$=41$, \code{max\_support}$=5$); 2-opt i polerowanie
\code{transpile} nie znalazły poprawy. Lokalne optimum jest dla tego
wejścia wyraźnie ciasne.

% ---------------------------------------------------------------------
\subsection{Wektor B -- \texttt{D\_Lambda\_23} dostarczony przez
            użytkownika}
\label{sec:vec-B}

\begin{lstlisting}
D_Lambda_23 = np.array([
    1,
    omega3_bar, omega3_bar, omega3_bar, omega3_bar,
    1, 1, 1,
    omega3, 1, omega3, 1, 1, omega3, 1, omega3,
    1, 1, 1, omega3, 1,
    omega3_bar, omega3, omega3_bar, omega3_bar,
    omega3, omega3, omega3_bar,
    1, 1,
    omega3, omega3, omega3_bar,
    1, 1, omega3_bar,
])
\end{lstlisting}

\begin{center}
\begin{tabular}{l r r r r}
\toprule
\textbf{wariant} & \textbf{depth} & \textbf{cx} & \textbf{rz} &
\textbf{niez.\ wyr.\ Walsha}\\
\midrule
naiwny baseline (\emph{don't-care} = 1, bez optymalizacji) & 107 & 60 & 53 & 53\\
\code{method="both"} (RS + multi-restart SA)               &  94 & 56 & 45 & 45\\
\code{method="best\_of"} (dodaje LP, greedy, 2-opt)        & \textbf{80} & \textbf{52} & 33 & \textbf{33}\\
phase-gadget po \code{transpile} (\code{best\_of})         &  98 & 92 & 33 & 33\\
phase-gadget raw (\code{best\_of})                          & 139 &130 & 33 & 33\\
\bottomrule
\end{tabular}
\end{center}

\noindent
\code{verify\_on\_code\_space} $\approx 1{,}4\times 10^{-15}$ dla
wszystkich trzech wariantów.

\paragraph{Rozkład „kto co znalazł” (przy domyślnym budżecie).}
\begin{lstlisting}
random_search             cost = 190.600  cx_est=186  terms=43  max_supp=6
sa(x6)                    cost = 186.750  cx_est=182  terms=45  max_supp=5
lp(continuous->alphabet)  cost = 264.600  cx_est=258  terms=63  max_supp=6
lp + greedy               cost = 192.650  cx_est=188  terms=44  max_supp=5
greedy(x6)                cost = 186.750  cx_est=182  terms=45  max_supp=5
two_opt(walsh)            cost = 133.550  cx_est=130  terms=33  max_supp=5  <-- zwyciezca
polish_target='cx_then_depth': brak poprawy
\end{lstlisting}

\noindent
Duży skok pochodzi od 2-opt: SA, greedy i zaokrąglony LP utykają
w okolicy $\sim$45 wyrazów Walsha; 2-opt znajduje basen 33-wyrazowy w
\emph{jednym} przebiegu. Łączna poprawa na tym wektorze:
\[
\begin{array}{rcccc}
\mathrm{depth} & : & 107 \to 80         & (-25\%) \\
\code{cx}      & : &  60 \to 52         & (-13\%) \\
\code{rz}      & : &  53 \to 33         & (-38\%) \\
\text{wyr.\ Walsha} & : & 53 \to 33     & (-38\%)
\end{array}
\]

% ---------------------------------------------------------------------
\subsection{Reprodukcja obu tabel wyników}

\begin{lstlisting}
import numpy as np
from d6_diagonal_optimizer import (
    optimize_d6_diagonal, verify_on_code_space, _example_d_lambda_2_3
)

omega3 = np.exp(2j * np.pi / 3); omega3_bar = np.conjugate(omega3)
D_Lambda_23 = np.array([
    1,
    omega3_bar, omega3_bar, omega3_bar, omega3_bar,
    1, 1, 1,
    omega3, 1, omega3, 1, 1, omega3, 1, omega3,
    1, 1, 1, omega3, 1,
    omega3_bar, omega3, omega3_bar, omega3_bar,
    omega3, omega3, omega3_bar,
    1, 1,
    omega3, omega3, omega3_bar,
    1, 1, omega3_bar,
])

for label, diag in [("vector A", _example_d_lambda_2_3()),
                    ("vector B", D_Lambda_23)]:
    res = optimize_d6_diagonal(
        diag,
        phase_alphabet="default",
        max_iters=4000,
        seed=2024,
        order="row-major",
        method="best_of",
        verbose=False,
    )
    m = res["metrics"]
    err = verify_on_code_space(res["baseline_circuit"], diag)
    print(f"{label:9s}  baseline depth={m['baseline_depth']:>3} "
          f"cx={m['baseline_cx']:>3} rz={m['baseline_rz']:>3}  "
          f"walsh_terms={res['num_nonzero_terms']:>2}  err={err:.1e}")
\end{lstlisting}

\noindent
Spodziewane wyjście:
\begin{lstlisting}
vector A   baseline depth= 85 cx= 50 rz= 41  walsh_terms=41  err=1.5e-15
vector B   baseline depth= 80 cx= 52 rz= 33  walsh_terms=33  err=1.4e-15
\end{lstlisting}

% =====================================================================
\section{Wskazówki praktyczne}
\label{sec:tips}
% =====================================================================

\begin{itemize}[leftmargin=*]
  \item Dla wejść o wartościach w $\w_{3}$ (AME(4,6) i pochodne) zostaw
        \code{phase\_alphabet="default"} = $\{0,\,2\pi/3,\,-2\pi/3\}$.
        Poszerzanie alfabetu zwykle \emph{wprowadza} dodatkową treść
        Walsha zamiast usuwać istniejącą.
  \item Dla trudniejszych przypadków zwiększ \code{max\_iters} oraz
        \code{n\_restarts}.
  \item Ustaw \code{polish\_target=None}, jeśli polerowanie
        \code{transpile} nie jest potrzebne (dodaje $\sim$1 sekundę;
        przydatne tylko, gdy zastępczy koszt Walsha jest luźny).
  \item Jeżeli chcesz obwód czysto \code{CX}+\code{RZ} (bez
        \code{transpile}), użyj \code{result["phase\_gadget\_circuit"]}.
        Pamiętaj, że dla obu wektorów przykładowych \emph{baseline}
        \code{DiagonalGate} jest \emph{tańszy} niż naiwny phase-gadget,
        bo Qiskit-owa synteza \code{DiagonalGate} stosuje
        Gray-code-owe scalanie, którego nasza prosta implementacja nie
        wykonuje.
\end{itemize}

% =====================================================================
\section{Otwarte pomysły / dalsze prace}
\label{sec:future}
% =====================================================================

\begin{enumerate}[leftmargin=*]
  \item \textbf{2-opt z metryką \code{transpile}.} Obecne polerowanie
        to 1-opt nad \code{transpile()}. Pełne 2-opt to
        $\binom{28}{2}\cdot 9 \approx 3{,}4$~tys.\ wywołań
        \code{transpile} na przebieg ($\sim$3~min); czasem wyciska
        kolejne kilka \code{cx}/głębokości. Można ograniczyć
        rozważane pary do tych z największym wkładem w koszt Walshowy.
  \item \textbf{Eksperymenty z szerszym alfabetem.} Spróbować
        $\Alphabet=\{0,\,\pm\pi/3,\,\pm 2\pi/3,\,\pi\}$ dla obu wektorów.
        Oczekiwanie: zazwyczaj \emph{gorsze} dla wejść o strukturze
        $\w_{3}$, ale potencjalnie lepsze dla wejść o drobniejszej
        strukturze fazowej.
  \item \textbf{Ciaśniejsza dolna granica.} Przy 28 swobodnych
        zmiennych rzeczywistych i 63 niezerowych współczynnikach
        Walsha naiwna dolna granica liczby niezerowych wyrazów
        wynosi $\sim 35$. Wynik 33 dla wektora B \emph{już to bije},
        więc macierz $B$ ma efektywny rząd $<28$ dla wejść
        $\w_{3}$ (znane zależności kolumnowe). Dokładny dyskretny
        optimum znalazłoby małe MIP (28 zmiennych ternarnych),
        rozwiązywane w sekundach przez CPLEX/Gurobi/SCIP.
  \item \textbf{LP $\to$ MIP, warm start.} Punkt ciągły LP można
        wykorzystać jako start dla MIP-a kodującego ograniczenia
        alfabetu.
  \item \textbf{Lepsza synteza phase gadgets.} Naiwna drabinka
        \code{CNOT} kosztuje $2(k-1)$ \code{CX} na wyraz. Synteza
        z grupowaniem Gray-code'owym lub \emph{ParityNetwork}
        (Amy/Maslov) istotnie zmniejszyłaby \code{cx} surowego
        obwodu phase-gadgetowego, prawdopodobnie do kilku procent
        powyżej baseline'u \code{DiagonalGate}.
\end{enumerate}

% =====================================================================
\section{Układ plików i status}
\label{sec:layout}
% =====================================================================

\begin{lstlisting}
QuditsOnQubits/QuditsOnQubits/
|-- d6_diagonal_optimizer.py                # modul
|-- docs/
|   |-- d6_diagonal_optimizer_README.md     # README po angielsku
|   |-- d6_diagonal_optimizer_README.tex    # ten dokument
|   |-- d6_diagonal_optimizer_notes.md      # snapshot notatek z sesji
`-- find_best_diagonal_decomposition.py     # pokrewny, ogolniejszy modul
                                            # (bez optymalizacji don't-care)
\end{lstlisting}

\noindent
Uruchomienie \code{python d6\_diagonal\_optimizer.py} wykonuje
self-testy oraz domyślny przykład. Lint przechodzi czysto, moduł
przechodzi self-testy, oba wektory przykładowe dają metryki podane w
sekcji~\ref{sec:results} na Qiskit 2.3.1, NumPy z \code{bitwise\_count}
(lub powolnym fallbackiem) oraz SciPy 1.x.

\bigskip\hrule\bigskip
\noindent
\textbf{Powiązany dokument źródłowy:} \\
\code{docs/d6\_diagonal\_optimizer\_README.md} (oryginał po angielsku).

\end{document}
